% Power-domain material for the consolidated "Reset, clocks and power"
% chapter. Input by TRM.template.tex, which owns the \chapter; this file
% carries the \section{Power domains} (label s:powerdomains, referenced by the
% template) and \subsection headings under it. Extracted from the PWRCTRL intro (2026-08); the
% PWRCTRL chapter keeps the register-level sequencer and boot-gate description.
% Hart numbering is derived: gateable harts are 1--\MaxHartIndex{} at every
% hart count, and \PwrGateMask is the PWRCR mask those bits make.
\section{Power domains} \label{s:powerdomains}

The chip has one always-on power domain and one switchable domain per channel tile. Everything the rest of the chip needs (the arbiter, the boot ROM, the shared RAM, the peripherals, the pad ring and hart 0) sits on the unswitched \texttt{VDD} rail. Each channel tile, harts 1 to \MaxHartIndex{}, has a rail of its own behind PMOS header switches, with isolation cells clamping every signal that leaves the tile while it is unpowered, and a per-tile hardware sequencer in \peripheral{PWRCTRL} that applies the controls in the legal order in both directions. Figure \ref{fig:pwr-domains} draws the domains and what one gate bit does.

% Generated by LatexUserGuide.GeneratePowerDomainDiagram (make chip).
\begin{figure}[htpb]
	\centering
	\resizebox{\linewidth}{!}{\input{include/PowerDomainDiagram.tex}}
	\caption{The chip's power domains and what a gate bit does.}
	\label{fig:pwr-domains}
\end{figure}

\iforchpresent
Hart 0, the orchestrator, is soft logic in the always-on centre band with no header switches of its own; it is permanently powered and cannot be gated.
\else
Hart 0, the management hart, is the same tile as the others but is wired always-on: its power-domain controls are tied inactive and the controller gives it no sequencer row, so it is permanently powered and cannot be gated.
\fi{}
All domains power up on at reset, so chip boot does not depend on the controller, and \peripheral{PWRCTRL} is inert until software gates a tile.

\subsection{Cold gating}

Gating is cold: a gated tile consumes no dynamic or switched-leakage power and retains nothing, including its register file, CSRs and private TCM. The isolation clamps drive every outbound tile signal to its inactive level, which is identical to the tile's reset state, so the rest of the chip observes a gated tile as a silent one; a read of its TCM aperture completes normally and returns zeros. When the tile is ungated it cold-boots through the shared boot ROM, parks, and is relaunched with the loader-row and software-interrupt sequence of Section \ref{ss:multicore-boot}. Any data a tile must keep across a power cycle is staged to shared RAM before gating.

\subsection{Memory block power}

The ROM, the two per-hart memory blocks and the first four shared RAM banks each have a power gate in \register{BLOCKPWR} (Section \ref{ss:system-blockpwr}). A gated block halves its leakage and loses its contents; an access to a gated block returns undefined data. Every block powers up on at reset.

\subsection{Hart sleep}

Each hart can stop its own clock with the \asminline{sleep} instruction and is woken by an interrupt (Section \ref{ss:sleep}). Sleep is a clock gate, not a power gate: the hart keeps its state and resumes where it stopped, so it is the right tool for a hart that waits for an event, while cold gating is the right tool for a tile that has no work for a long time.

\subsection{Boot gate and field power}

In configurations wired for field power (the harvested-energy NFC operating mode), \peripheral{PWRCTRL} also decides when the chip may start executing, so that the harts are held quiescent until the harvested supply can carry a transaction. Two package pins take part: a power-good input (\texttt{PGOOD}, pin 69, P6.7) driven by an external supply supervisor, and a boot-mode strap (pin 68, P6.6) that selects harvested boot. Both pads reset to inputs with pull-downs, so an unconnected pin reads as power-not-good and normal boot, and a board that does not use the feature behaves as if it did not exist.

The boot gate is a hold-in-reset: an internal release signal is combined into every hart's outer reset, hart 0 included, so a held hart fetches nothing and issues no bus request, and the rest of the chip cannot distinguish it from a hart still in power-on reset. The strap is sampled once, eight MCLK cycles after reset release, and sets the hardware defaults of the gate: a strapped board arms the gate, makes \texttt{PGOOD} a release condition and re-holds on brownout with no firmware at all, which resolves the problem that the gate holds the very harts that would otherwise program it. Software may add overrides through \register{PWRWAKE}: arm the gate, add \texttt{PGOOD} or the NFC field-detect level as release conditions, force a release, and choose between a one-shot release and re-holding whenever a release condition later drops. Field detect used this way is a wake path that consumes no interrupt vector; the same level remains available as an ordinary interrupt to firmware that is already running.

With re-hold enabled, a \texttt{PGOOD} drop re-asserts the hold on the next cycle and its return releases the harts into a fresh boot; since the chip keeps no retention, a full cold boot is the accurate model of losing and regaining the harvested rail. Two schemes are supported on the same silicon. In the batteryless scheme the gate is the mechanism: the chip holds until \texttt{PGOOD}, cold-boots the harvested-boot ROM path (Section \ref{ss:multicore-harvested}), serves the reader, and re-holds when the field collapses. In the supercap duty-cycled scheme an external buffer stays charged, the gate is left disarmed after the first boot, and the service loop sleeps between transactions and resumes on the field-detect interrupt without re-entering reset. The register-level mechanism and its timing are in the \peripheral{PWRCTRL} chapter.
